All studies in this report refer to a single reference configuration.
It is described here once; later chapters cite this chapter rather than
restating it.

\paragraph{Architecture.} The model follows the OpenMythos design
\citep{gomez2024openmythos}: 43,974,722 parameters, model dimension 512,
6 unshared prelude layers and 2 unshared coda layers framing a shared
recurrent block executed for 4 loops, 8 experts per mixture-of-experts
feedforward layer, grouped-query attention, and a GPT-2 byte-pair
vocabulary of 50,257 tokens. The prelude/coda structure is a descendant
of the Universal Transformer \citep{dehghani2019universal} family, with
parameter sharing in the spirit of ALBERT \citep{lan2020albert}.
Adaptive-depth variants allocate iteration per token rather than fixing
it globally \citep{elbayad2020depth,raposo2024mixture}; the reference
model here fixes the loop count at training time, which is what makes
the test-time extrapolation question sharp.

\paragraph{Training.} The reference run trains on TinyStories
\citep{eldan2023tinystories} for 4,850 steps (120,004,608 tokens,
approximately 152 minutes on a single Kaggle T4 GPU in fp16). Training
loss falls from 9.5242 to 1.5263. Qualitative inspection shows a clear
progression from incoherent output early in training to fluent
multi-sentence stories by the end.

\paragraph{Attribution.} The OpenMythos architecture is due to Gomez
\citep{gomez2024openmythos} (GitHub repository, 2024). The 44M training
run, the loop-depth sweep, and all analyses in this report are the
author's own work; no released checkpoints are used.
